%%=============================================================================
%% Voorwoord
%%=============================================================================

\chapter*{\IfLanguageName{dutch}{Woord vooraf}{Preface}}%
\label{ch:voorwoord}

%% TODO:
%% Het voorwoord is het enige deel van de bachelorproef waar je vanuit je
%% eigen standpunt (``ik-vorm'') mag schrijven. Je kan hier bv. motiveren
%% waarom jij het onderwerp wil bespreken.
%% Vergeet ook niet te bedanken wie je geholpen/gesteund/... heeft

Tijdens mijn studie aan de HoGent kreeg ik een passie voor netwerken en het opzetten van servers. Ik heb ongelooflijk veel tijd besteed aan het opzetten en troubleshooten van mijn eigen toestellen. Toen ik een topic moest kiezen voor mijn bachelorproef, vond ik het best lastig om er een te vinden waar ik motivatie van kreeg. Uiteindelijk besefte ik dat ik al een ideale testomgeving had die in mijn kamer stond.
Ik besloot om de uitdaging aan te gaan, en mijn eigen netwerk te gebruiken als een academische Proof-of-Concept. Wat begon als een hobby, zoals het installeren van Proxmox voor virtuele machine op een oude computer, werd al snel complexer door het toevoegen van een Mikrotik Router en een managed switch.
Het combineren van mijn \textit{homelab} met een gestructureerde academische methodologie was een interessante rollercoaster. Ik moest mijn kennis vergroten en leren hoe protocollen zoals MQTT en Zigbee precies werken, en hoe ik al die apparaten op mijn netwerk kon aansluiten. Bij het leren van hoe Suricata en Zeek kunnen ingesteld worden als IDS, viel Suricata best mee. Maar het zorgen dat Zeek dezelfde meldingen kon maken zoals Suricata was wat lastiger.

Voordat ik dit stuk afsluit, wil ik mijn promotor, Lieven Smits, en mijn co-promotor, Martijn Saelens, heel erg bedanken voor hun hulp en geduld. Als ik vastliep waren jullie er altijd om me verder te helpen met advies. Tot slot wil ik ook HoGent bedanken voor de begeleiding en de ondersteuning tijdens dit traject.
